\sscp{Complexity in numerals}{02-03-01-04}
It has been observed that there are a number of Austronesian
languages in Wallacea (such as in some languages of Flores,
Timor, \ili{Aru}, \ili{Buru}) and out into the Pacific that have complexity
in numerals between 6
and 10, in contrast with the simple numerals inherited from PMP,
and that similar patterns are found in Papuan languages in the
region, such as in \tsc{\ili{Alor}-Pantar} languages.
\citep[111f]{SchapperHammarstrom2013,Schapper2015}.
Some complexity in numerals diverging from the PMP forms is also
found west of Wallacea,
such as in \ili{Malay} \it{dəlapan} ‘eight’
and \it{səmbilan} ‘nine’, and the subtractive strategy of
\tsc{South Sulawesi} languages ‘minus-two = 8’
and ‘minus-one = 9’ (see \citealp[130--133]{GrimesGrimes1987}).
\trf{02-03} shows the simple PMP forms,
along with several languages in Wallacea that show retentions
of the PMP numerals below 10. (\it{m.} = metathesis)

\begin{table}[h]
	\centering\tcp{Retention of simple numerals inherited from PMP}{02-03}
	\begin{threeparttable}
	\st{2pt}\begin{tabular}{lllllll}\lsptoprule
	& Subgroup	& BL & TB & CT & AS & STB \\
	& PMP	& \ili{Dhao} & \ili{Tii} & \ili{Kemak} & \ili{Yalahatan} & \ili{Yamdena} \\ \midrule
1	&	*[ə/i/a]sa	&	əʧi \it{(m.)}	&	esa	&	sea	&	isa, sa	&	sa, \tcg{(lese)}		\\
2	&	*duha	&	ɗua	&	dua	&	rua	&	rua	&	du, dua, due		\\
3	&	*təlu	&	təlu	&	telu	&	telu	&	telu	&	tel, tely		\\
4	&	*əpat	&	əpa	&	haa \it{(m.)}	&	paat \it{(m.)}	&	ata	&	fat		\\
5	&	*lima	&	ləmi \it{(m.)}	&	lima	&	lima	&	lima	&	lim		\\
6	&	*ənəm	&	əna	&	nee  \it{(m.)}	&	neem \it{(m.)}	&	nome	&	nem		\\
7	&	*pitu	&	pi{\DZ}u	&	hitu	&	iʧu	&	itu	&	itw		\\
8	&	*walu	&	aru	&	falu	&	balu	&	walu	&	walu, walw		\\
9	&	*siwa	&	ʧeo	&	sio	&	sibe	&	siwa	&	siu, siw		\\
10	&	*sa-ŋa-puluq	&	ʧa-ŋuru	&	salahunu \it{(m.)}	&	gulu\su{†} (‘tens’)	&	\tcg{(hutu sa)}	&	\tcg{(buty)}		\\
		\lspbottomrule
	\end{tabular}
			\begin{tablenotes}\tnsize
				\item[†]	{\ili{Kemak} has \it{sapulu} ‘ten’
									(possibly a \ili{Malay} loan, but would also
									be a regular inheritance from *sa-puluq).
									Multiples of ten use \it{gulu} as in \it{gulu rua} ‘20’.}
			\end{tablenotes}
	\end{threeparttable}
\end{table}

Other patterns are also found in the region.
These include additive patterns often built around four,
five, or six plus a lower number,
and subtractive patterns, such as those mentioned for
\tsc{South Sulawesi} languages above.
These complex patterns are illustrated in \trf{02-04}
with complex numerals shaded.
%The \ili{Rongga} data below are adapted from \citet[113]{Schapper2015}.

\begin{table}[h]
	\centering\tcp{Complexity in numerals between 6--10 (Austronesian)}{02-04}
\begin{threeparttable}
	\st{4.5pt}\begin{tabular}{lll<{\hspace{-\tabcolsep}}>{\hspace{-\tabcolsep}\,}ll<{\hspace{-\tabcolsep}}>{\hspace{-\tabcolsep}\,}ll<{\hspace{-\tabcolsep}}>{\hspace{-\tabcolsep}\,}l}
\lsptoprule
	& 	 & BL && CT && SB &\\
	& PMP & \ili{Rongga} && S. \ili{Mambae} && \ili{Buru} &\\ \midrule
1	&	*[ə/i/a]sa	&	esa, sa	&		&	\tcg{(ida, iid, kiid)}	&		&	\mc{2}{l}{em-sia-n, saa \it{(m.)}}				\\
2	&	*duha	&	ɹua	&		&	ruu, rua	&		&	rua	&			\\
3	&	*təlu	&	telu	&		&	teul, telu	&		&	telo	&			\\
4	&	*əpat	&	\tcg{(wutu)}	&		&	faat, fata	&		&	paa \it{(m.)}	&			\\
5	&	*lima	&	lima	&		&	liim, lima	&		&	lima	&			\\
6	&	*ənəm	&	lima esa 	{\cgr}&	{\cgr}	[5+1]&	liim nai-nida 	{\cgr}&	[5+1]	{\cgr}&	nee \it{(m.)}	&			\\
7	&	*pitu	&	lima ɹua 	{\cgr}&	{\cgr}	[5+2]&	liim nai-rua {\cgr}	&	[5+2]	{\cgr}&	pito	&			\\
8	&	*walu	&	ɹua mbutu 	{\cgr}&	{\cgr}	[2×4]&	liim nai-telu	{\cgr}&	[5+3]	{\cgr}&	et-rua	{\cgr}&	{\cgr}	[-2]	\\
9	&	*siwa	&	tara esa 	{\cgr}&	{\cgr}	[-1]&	liim nai-fata	{\cgr}&	[5+4]	{\cgr}&	eʧia /et-sia/\su{†}	{\cgr}&	{\cgr}	[-1]	\\
10	&	*sa-ŋa-puluq	&	sambulu	&		&	saguul	&		&	polo	&			\\
12	&		&		&		&	\mc{2}{l}{saguul resi ruu}			&	\mc{2}{l}{polo geran rua}				\\
20	&		&		&		&	rua-nulu	&		&	po-rua	&			\\
80	&		&		&		&	ual-nulu\su{‡}	&		&	polo t-rua	&			\\
		\lspbottomrule
	\end{tabular}
			\begin{tablenotes}\tnsize
				\item[†]	{\ili{Buru} retains *w = \it{w}, so the form for ‘nine’
									is not formed through a loss of *w from *siwa.}
				\item[‡]	{Reflexes of the PMP forms occur
									in multiples of ten in South \ili{Mambae},
									e.g. \it{ual-nulu} ‘80’}
			\end{tablenotes}
	\end{threeparttable}
\end{table}

In \trf{02-04} one can see that the \ili{Mambae} system is built on
adding numerals 1-4 onto a quinary base of five.
Note that \ili{Mambae} (in \trf{02-04},
with its innovative system,
and \ili{Kemak} (in \trf{02-03}),
with its conservative retention of the PMP system,
are neighbouring languages both within the \tsc{\ili{Central} Timor} subgroup.
The \ili{Buru} system of minus-two
and minus-one is discussed in more detail in \srf{08-04}.
The \ili{Rongga} system is a hybrid,
adding numerals 1--2 onto a base of five,
but also doing different things for 8 and 9.

Note also that the \ili{Rongga} system diverges from the PMP system
with its numerals 4
and 8 (\ili{Rongga} \it{wutu} ‘4’ is likely from PCEMP **butu ‘group,
crowd, flock, bunch, cluster’).
Looking at ethno-mathematics in \ili{Amarasi},
\citet{GrimesBani2011} show that a number of terms for counting
large quantities of corn are based on how they are tied,
and revolve around bases of 4
and 8, or 5 and 10.
For example a \it{rean} is 400 ears of corn when tied with a
base of 5 ears,
whereas a \it{rean} is 384 ears of corn when tied with a base of 4 ears.
Similarly \ili{Dela} terms \it{lombeʔ,
taliʔ}, and \it{ritiʔ} have variable readings according to how
many ears of corn are tied in a bundle \citep[116]{Tamelan2021}.

Similar additive systems on a base of 5,
and subtractive systems on a base of ten are found in Papuan
languages in the region,
such as \ili{Alor} languages \ili{Klon}
and \ili{Wersing}, shown in \trf{02-05}
with complex forms shaded.

\begin{table}[h]
	\centering\tcp{Complexity in numerals between 6--10 (Papuan)}{02-05}
\begin{threeparttable}
	\st{6pt}\begin{tabular}{lllll}\lsptoprule
	& \mc{2}{l}{TAP, \tsc{\ili{Alor}-Pantar}}& \mc{2}{l}{TAP, \tsc{\ili{Alor}-Pantar}}\\
	& \mc{2}{l}{Wersing}& \mc{2}{l}{\ili{Klon} }\\ \midrule
1	&	no	&		&	nuk	&			\\
2	&	yoku	&		&	ʊrʊk	&			\\
3	&	tu	&		&	toŋ	&			\\
4	&	arasoku	&		&	ut	&			\\
5	&	wetiŋ	&		&	ɪwɪh	&			\\
6	&	wetinuŋ, tlaam	{\cgr}&	[5+1], [6]	{\cgr}&	tlaan	&			\\
7	&	wetiŋ yoku	{\cgr}&	[5+2]	{\cgr}&	usoŋ	&			\\
8	&	wetiŋ tu	{\cgr}&	[5+3]	{\cgr}&	tid-ʊrʊk	{\cgr}&	[-2]	{\cgr}	\\
9	&	weit arasoku\su{†}	{\cgr}&	[5+4]	{\cgr}&	tukeɲ-nuk	{\cgr}&	[-1]	{\cgr}	\\
10	&	ada-yoku	&		&	kar nuk	& ‘10’			\\
12	&	\mc{2}{l}{ada-yoku gresiŋ yoku}			&	kar nuk awaa ʊrʊk	& ‘12’			\\
20	&	ada-yoku mi yoku	&		&	kar ʊrʊk 	& ‘20’			\\
		\lspbottomrule
	\end{tabular}
			\begin{tablenotes}\tnsize
				\item[†]	{
									The first part of \it{weit arasoku}
									shows regular final CV → VC metathesis
									before vowel-initial words.}
			\end{tablenotes}
	\end{threeparttable}
\end{table}

The \ili{Wersing} (Papuan) system parallels the additive quinary system
found in South \ili{Mambae} (Austronesian),
adding numerals 1-4 onto a base of 5.
It is likely that the additive quinary system found in Austronesian
languages is due to contact with Papuan languages.

The \ili{Klon} system parallels the subtractive system found in \ili{Buru} and elsewhere.
However, since the minus-two, minus-one strategy is also found
outside of Wallacea in west Indonesia
and \tsc{South Sulawesi} languages,
perhaps the best one can say is that languages within Wallacea
that use this strategy have had pressure from contact with other
languages that also do that,
but that includes other Austronesian languages as well as Papuan languages.

Neighbouring languages in the same node of the same subgroup
sometimes do different things, with one language retaining the Austronesian forms,
and its neighbour doing something else, such as Northwest \ili{Mambae} and
South \ili{Mambae} having complex numerals for 6--9,
but \ili{Central} \ili{Mambae} retaining the Austronesian forms
(\trf{05-03}, \prf{tab:05-03}).
%See also \srf{08-04} for discussion of numerals in \tsc{Taliabu}
%and \tsc{Sula-Buru}.

\frf{02-17} shows the distribution of complex numerals in our
target region, as well as neighbouring areas.
Note that \frf{02-17} does not distinguish between the different
strategies for forming complex numerals,
such as additive, subtractive and so on.
Note that with only three known exceptions (Proto-Malayic,
Enggano, and \ili{Ilongot}) complex numerals in Malayo-Polynesian languages
are only found in Sulawesi
and \ili{east} of the Blust line.\footnote{
		Several Austronesian languages
		of Taiwan also have complex numerals.}
Not all instances of
complex numerals are necessarily due to contact with Papuan languages.
The strongest evidence for such contact comes from non-decimal
bases, while additive/subtractive numerals for 8
and 9 are less likely due to Papuan contact.
See \citet{SchapperHammarstrom2013}
and \citet[111--114]{Schapper2015} for more discussion.

\begin{figure}
	\includegraphics[width=\columnwidth]{numerals.pdf}
	\fcp{Complex numerals in eastern Indonesia and Timor-Leste}{02-17}
\end{figure}

Looking briefly at order in the number NP,
the pattern of NUM + NOUN commonly found in Austronesian languages
west of the Blust line is preserved in western
and southern parts of the \tsc{\ili{Bima}-Lembata} subgroup,
with \tsc{\ili{Havu}-Dhao} having optional ordering.
The \tsc{Flores-Lembata} node has shifted to NOUN + NUM,
which is also the dominant order found in Papuan languages in
the region \citep[357--365]{Fricke2019}.
Other Wallacean subgroups also tend to have NOUN + NUM,
mirroring the order in Papuan SOV languages.